\chapter{Resultados}

Este capítulo apresenta os resultados obtidos na avaliação da ferramenta VANSI. A análise integra dados quantitativos de interação coletados automaticamente pela ferramenta com respostas qualitativas dos questionários de \textit{feedback}, organizando-se em cinco seções principais: caracterização dos dados coletados, análise dos padrões de interação, avaliação de usabilidade, eficácia pedagógica e avaliação geral.

\section{Caracterização dos dados coletados}
\label{sec:caracterizacao-participantes}
Os 23 participantes da avaliação apresentaram o perfil descrito na Tabela \ref{tab:perfil-participantes}. Nota-se que a maioria (60,9\%) não possuía experiência prévia com algoritmos de análise sintática, configurando uma amostra representativa do público-alvo da ferramenta.

\begin{mytabela}{tab:perfil-participantes}{Perfil dos participantes da avaliação (N=23)}{X[l]lcc}
	\hline
	\Block{1-2}{\textbf{Característica}} & & \textbf{Quantidade} & \textbf{Percentual} \\
	\hline
	\Block{2-1}{Experiência prévia com \textit{parsing}} &  Sim & 9 & 39,1\% \\
	& Não & 14 & 60,9\% \\
	\addlinespace

	\Block{3-1}{Nível de conhecimento autorrelatado} & Iniciante & 8 & 34,8\% \\
	& Intermediário & 12 & 52,2\% \\
	& Avançado & 3 & 13,0\% \\
	\addlinespace

	\Block{2-1}{Uso prévio da ferramenta} & Sim & 0 & 0\% \\
	& Não & 23 & 100\% \\
	\hline
\end{mytabela}

Durante as 7,7 horas totais de atividade, o sistema de \textit{log} da VANSI registrou 88.301 eventos de interação. Com 23 estudantes participantes, a média de 3.839 eventos por estudante indica alto engajamento durante a sessão de avaliação. O tempo médio de uso foi de 20,1 minutos. Estudantes sem experiência prévia gastaram em média 28,5 minutos, enquanto aqueles com experiência utilizaram 15,8 minutos, sugerindo que a familiaridade com os conceitos reduziu o tempo necessário para execução das sete tarefas.

A Tabela \ref{tab:volume-dados} detalha o volume de dados coletados por categoria de evento. Os eventos de \textit{mouse} representaram a maioria dos registros, seguidos pelas interações com controles de fluxo e trocas de aba.

\begin{mytabela}{tab:volume-dados}{Volume de dados coletados durante a avaliação}{Xlr}
	\hline
	\textbf{Tipo de Evento} & \textbf{Total} & \textbf{Média por Estudante} \\
	\hline
	Eventos de mouse (posição + clique) & 87.421 & 3.801                    \\
	Trocas de aba & 412 & 17,9                                              \\
	Manipulações de janelas flutuantes & 147 & 6,4                          \\
	Controles de fluxo utilizados & 321 & 14,0                              \\
	\hline
	\textbf{Total de eventos} & \textbf{88.301} & \textbf{3.839}            \\
	\hline
\end{mytabela}

\section{Análise dos padrões de interação}
\label{sec:padroes-interacao}

Nesta seção são examinados os comportamentos de navegação e interação dos participantes com a ferramenta, com base nos registros automáticos de eventos. Os resultados estão organizados em duas subseções, distribuição do tempo por aba e padrões de navegação e interação.


\subsection{Distribuição do tempo por aba}

A distribuição do tempo gasto em cada aba da ferramenta revelou padrões consistentes entre os participantes. A aba \gls{slr} recebeu a maior atenção com 31,7\% do tempo total, seguida pela aba \gls{ll}(1) com 25,4\% e \gls{lr}(1) com 20,1\%. A aba de Gramática, destinada à entrada inicial dos dados, consumiu 15,5\% do tempo, enquanto a aba de Análise de \textit{String}, utilizada principalmente para validação final, representou 7,3\% do tempo total. Esta distribuição sugere que os algoritmos \gls{slr} e \gls{ll}(1), tradicionalmente mais desafiadores para os estudantes, demandaram maior tempo de exploração e compreensão. A Figura \ref{fig:timeline} mostra a linha de tempo de um dos \textit{logs}.

\begin{mypicture}
	{\label{fig:timeline}}
	{Linha de tempo de uma sessão de uso da ferramenta}
	{figuras/timeline.png}
	{fornecido pelo próprio autor}
\end{mypicture}


\subsection{Padrões de navegação e interação}

A análise dos \textit{heatmaps} de cliques revelou padrões de interação distintos entre as abas. Na aba Tarefa, observou-se distribuição esparsa das interações pelos campos de entrada. Nas abas algorítmicas, porém, identificou-se concentração em elementos visuais específicos, a tabela de \textit{parsing} na \gls{ll}(1), o autômato \gls{lr}(0) no \gls{slr}, e a árvore sintática na análise de \textit{strings}.

A Figura \ref{fig:heatmap} apresenta um \textit{heatmap} de uma sessão de avaliação na aba SLR. Esse \textit{heatmap} mostra a concentração de interações em elementos-chave da interface e padrão de uso esperado da ferramenta. Ainda que a maioria dos elementos tenha sido explorada, algumas funcionalidades como o botão ``Ir'' e a janela de pseudocódigo não receberam interações. Esse padrão de uso pode ser explicado pela quantidade grande de funcionalidades disponíveis e a falta de familiaridade com a ferramenta, sendo necessário um período de adaptação para explorar todas as funcionalidades. Também é possível que a atividade proposta não tenha demandado o uso de todas as funcionalidades, o que é um aspecto a ser considerado no \textit{design} de tarefas para futuras avaliações.

\begin{mypicture}	{\label{fig:heatmap}}
	{\textit{Heatmap} de cliques na aba Tarefa}
	{figuras/heatmap.png}
	{fornecido pelo próprio autor}
\end{mypicture}

Os controles de fluxo foram utilizados extensivamente, com média de 14 ações por estudante. O botão ``Avançar'' representou 45,2\% das ações, seguido por ``Voltar'' com 32,4\% e ``Pular para passo'' com 12,1\%. A predominância do uso sequencial (``Avançar'' e ``Voltar'') sobre o acesso direto (``Pular para passo'') sugere que os estudantes preferiram percorrer os algoritmos de forma linear, revisando passos quando necessário, em vez de acessar etapas diretamente.

\section{Avaliação de usabilidade}

As respostas às questões de usabilidade apresentaram avaliação geral positiva. A Tabela \ref{tab:avaliacao-usabilidade} sumariza os resultados obtidos nas sete questões relacionadas à interface e usabilidade. A clareza dos elementos da interface obteve a maior nota média (4,3/5), com 91,0\% das respostas nas categorias 4 ou 5. O diálogo para copiar resultados também recebeu avaliação elevada (4,5/5), indicando que a funcionalidade de exportação foi bem recebida. A facilidade de navegação obteve média de 4,1/5, com 85,0\% de avaliações positivas.

\begin{mytabela}{tab:avaliacao-usabilidade}{Avaliação de usabilidade da ferramenta (escala 1-5)}{Xccc}
	\hline
	\textbf{Aspecto Avaliado} & \textbf{Média} & \textbf{DP} & \textbf{\% 4-5} \\
	\hline
	Facilidade de navegação & 4,1 & ±0,8 & 85\% \\
	Clareza dos elementos & 4,3 & ±0,7 & 91\% \\
	Utilidade dos \textit{tooltips} & 4,0 & ±1,0 & 78\% \\
	Mover janelas flutuantes & 4,4 & ±0,6 & 92\% \\
	Minimizar/abrir elementos & 4,2 & ±0,9 & 83\% \\
	Organização das abas & 4,1 & ±0,8 & 87\% \\
	Diálogo para copiar resultados & 4,5 & ±0,5 & 96\% \\

	\textbf{Média Geral} & \textbf{4,2} & \textbf{±0,7} & \textbf{87\%} \\
	\hline
\end{mytabela}

A análise qualitativa dos \textit{logs} permitiu identificar dois padrões comportamentais associados a dificuldades de usabilidade. Sete estudantes (30,4\%) apresentaram um padrão de busca, caracterizado por múltiplas trocas rápidas de aba nos primeiros 5 minutos de uso, reduzindo progressivamente após período de adaptação. Nove estudantes (39,1\%) apresentaram uso sub-ótimo de controles, utilizando predominantemente o botão ``Avançar'' sem explorar funcionalidades de pular para passo.

\section{Eficácia pedagógica}

Os recursos pedagógicos da VANSI receberam avaliação consistentemente alta. A Tabela \ref{tab:avaliacao-pedagogica} apresenta os resultados das sete questões relacionadas aos recursos pedagógicos da ferramenta. A visualização passo a passo obteve a maior nota média (4,7/5) e o maior consenso entre os participantes, com 100\% das respostas concentradas nas categorias 4 e 5. Os controles de navegação (4,6/5) e a árvore sintática (4,5/5) também figuram entre os recursos mais bem avaliados, ambos com 96\% e 91\% de respostas positivas, respectivamente.

As animações e transições, implementadas para destacar mudanças nas estruturas de dados, receberam média de 4,4/5, com 91\% das avaliações entre 4 e 5. O pseudocódigo com \textit{breakpoints}, embora avaliado positivamente (4,2/5), obteve a menor média e o menor percentual de aprovação entre os recursos pedagógicos (83\%), possivelmente por demandar familiaridade prévia com a ferramenta. De modo geral, 91\% das respostas do questionário indicaram satisfação com os recursos pedagógicos oferecidos.

\begin{mytabela}{tab:avaliacao-pedagogica}{Avaliação dos recursos pedagógicos (escala 1-5)}{Xccc}
	\hline
	\textbf{Recurso Pedagógico} & \textbf{Média} & \textbf{DP} & \textbf{\% 4-5} \\
	\hline
	Visualização passo a passo & 4,7 & ±0,5 & 100\%\\
	Animações e transições & 4,4 & ±0,7 & 91\%\\
	Visualização do autômato & 4,3 & ±0,8 & 83\%\\
	Árvore sintática & 4,5 & ±0,6 & 91\%\\
	Controles de avançar/voltar & 4,6 & ±0,5 & 96\%\\
	Pseudocódigo com \textit{breakpoints} & 4,2 & ±0,9 & 83\%\\
	Diálogos informativos & 4,3 & ±0,7 & 87\%\\

	\textbf{Média Geral} & \textbf{4,4} & \textbf{±0,7} & \textbf{91\%}\\
	\hline
\end{mytabela}


\section{Avaliação geral}

A avaliação geral da ferramenta foi predominantemente positiva, com 86,9\% dos participantes classificando-a como ``Excelente'' e 13,1\% como ``Boa''. Nenhum participante avaliou a ferramenta como ``Ruim'' ou ``Péssima''. A consistência entre os diferentes tópicos avaliadas sugere que a VANSI foi bem recebida tanto em aspectos técnicos quanto pedagógicos. As questões sobre intenção de uso futuro revelaram alta aceitação da ferramenta. 90,0\% dos participantes responderam que usariam a ferramenta novamente para estudar algoritmos de compiladores, enquanto 85,0\% afirmaram que a recomendariam para outros alunos.

\subsection{Sugestões e comentários}

A análise temática dos comentários abertos identificou quatro categorias principais de \textit{feedback}, conforme apresentado na Tabela \ref{tab:temas-comentarios}. Dois participantes elogiaram a interatividade da ferramenta, destacando que poder interagir com cada passo tornou a aprendizagem muito mais ativa e compreensível.

\begin{mytabela}{tab:temas-comentarios}{Temas identificados nos comentários abertos}{Xc}
	\hline
	\textbf{Tema} & \textbf{Frequência} \\
	\hline
	Elogios à interatividade & 2 \\
	Sugestões para experiência mobile & 5 \\
	Sugestões de melhorias na interface & 3 \\
	\hline
\end{mytabela}

Dentre as sugestões específicas, dois aspectos foram particularmente mencionados pelos participantes. Cinco usuários que acessaram a ferramenta por dispositivos móveis apontaram a necessidade de melhorias na experiência mobile, com comentários como ``seria bom melhorar um pouco a visualização dos elementos pelo celular'' e ``no celular alguns botões ficam muito pequenos para clicar''.

Três estudantes sugeriram melhorias específicas na interface visual, com ênfase na clareza dos estados dos elementos interativos. Os comentários incluíam observações como ``as cores dos botões ativos/inativos poderiam ser mais diferentes'' e ``fica confuso saber quando um botão está desabilitado''. Essas observações apontam para oportunidades de aprimoramento de \textit{affordance} visual dos controles da interface, isto é, para melhorar a percepção de funcionalidade dos elementos da interface.
